% ---------------- title page ----------------
\begin{titlepage}
\thispagestyle{empty}
\centering
\vspace*{1.05in}

{\sffamily\fontsize{8.5}{10}\selectfont\color{emerald}\MakeUppercase{The Isegoria Mathematics Series \quad$\cdot$\quad Volume II}}

\vspace{0.42in}

{\sffamily\bfseries\fontsize{40}{44}\selectfont\color{navy} Structure\\[2pt] Before Algebra}

\vspace{0.26in}

{\color{emerald}\rule{2.7in}{0.9pt}}

\vspace{0.26in}

{\fontsize{13.5}{17}\selectfont\color{ink!85}
A complete textbook of solution technique}

\vspace{0.14in}

{\sffamily\fontsize{10}{14}\selectfont\color{ink!62}
Integrals \quad$\cdot$\quad Derivatives \quad$\cdot$\quad Differentials \quad$\cdot$\quad Limits \quad$\cdot$\quad Logic}

\vfill

\begin{minipage}{4.35in}
\centering\itshape\fontsize{10.2}{14}\selectfont\color{ink!78}
A hard problem is hard because of what you fail to notice in the first thirty
seconds, not because of the algebra that follows.
\end{minipage}

\vfill

{\sffamily\fontsize{9}{12}\selectfont\color{ink!58}
\the\year \quad$\cdot$\quad isegoria.io}

\vspace{0.5in}
\end{titlepage}

% ---------------- colophon ----------------
\thispagestyle{empty}
\vspace*{\fill}
\begin{minipage}{\textwidth}
\small\color{ink!72}
\textbf{Structure Before Algebra}\\
A complete textbook of solution technique in calculus and logic.\\[6pt]
Volume II of the Isegoria mathematics series. Volume I is
\emph{The Complete Daily Integral Field Guide}.\\[6pt]
This book was assembled from a working corpus of solution archetypes: the
ninety-six core generator families and the several hundred technique variants
catalogued from daily problem sets, deduplicated into the technique inventory
that organises these chapters. Every closed form, limit, derivative, solution
and count asserted in the text was verified numerically or by exhaustive
enumeration before it was typeset, and a number of errors in the source
catalogue were corrected in the process.\\[6pt]
Set in TeX Gyre Pagella with TeX Gyre Pagella Math and TeX Gyre Heros.
Typeset with XeLaTeX.\\[6pt]
Published at \texttt{isegoria.io}.
\end{minipage}
\vspace*{\fill}
\clearpage

% ---------------- how to use ----------------
\thispagestyle{plain}
\section*{\sffamily\bfseries\color{navy}\fontsize{20}{24}\selectfont How to use this book}
\vspace{4pt}
{\color{rulegray}\rule{\textwidth}{0.6pt}}
\vspace{8pt}

This is a book about recognition, not about theory. It assumes you can
differentiate, integrate, and manipulate an inequality, and it spends none of
its pages proving that you can. What it does instead is catalogue, as
completely as we could manage, the moves that actually solve problems, and
attach to each one the surface feature that should make you reach for it.

Every technique in the book appears as a numbered block with the same five
parts. The \textsc{cue} is the visual or structural trigger: what you see on
the page that should make this technique come to mind. The \textsc{method} is
the mechanism, stated in as few sentences as it takes. The worked example is
one full instance, chosen to be hard enough to be worth the space. The trap
box names the specific way that technique is most often botched, which is
usually more valuable than the method itself. And the see-also line names
siblings, because techniques come in families and knowing the family is most
of knowing when to switch.

Chapter numbers and technique numbers are stable: technique $22.7$ is the
seventh technique of Chapter 22, and it is referred to that way in the index.
Every chapter closes with a recognition matrix, a compressed table of surface
features against first moves, and Chapter 53 collects the most useful rows of
all of them into five cross-domain tables. If you own the book for one thing,
own it for those tables.

There are three reasonable ways to read it.

\textbf{Front to back}, which works: the parts are ordered so that each one
uses the previous. Limits come before derivatives because differentiability is
a limit; derivatives come before integrals because half of integration is
recognising a derivative; differential equations come after both; and the
logic part stands alone but shares the book's organising idea, which is that
difficulty is camouflage.

\textbf{By domain}, treating each part as a self-contained manual. Nothing in
the integral chapters depends on the logic chapters, and the cross-references
are by name rather than by number so that they survive being read out of
order.

\textbf{By cue}, which is how the book is designed to be used once you have
read it once. You have a problem; you look at the recognition matrices; you
find the row that matches the surface; you go to the chapter that owns the
move. The index at the back is by technique name for the same purpose.

A word on what the book leaves out. There are no exercises. This is deliberate:
problems are abundant and technique catalogues are not, and a reader who wants
practice has a daily supply of it. There are also no proofs of the standard
theorems, only of the facts a solver actually uses, and where a result is
quoted rather than derived the text says so.

\clearpage
